%
% Documento: Resumo
%

\begin{resumo}

Este trabalho apresenta o desenvolvimento de um protótipo de leitor RFID UHF passivo destinado ao inventário de bens patrimoniais. O recorte adotado concentra-se na camada física e embarcada do sistema, formada por um módulo YPD-R200, uma antena externa de 5~dBi, etiquetas passivas e um Arduino. A proposta substitui uma concepção inicial baseada em NFC, cuja curta distância de operação exige a aproximação individual do leitor a cada item e limita a conferência de ambientes com muitos bens.

O protótipo utiliza a interface UART TTL do YPD-R200. A biblioteca de comunicação foi adaptada para o Arduino Uno/Nano por meio da abstração \texttt{Stream} e da biblioteca \texttt{SoftwareSerial}. Foram implementados o envio de comandos, a configuração da potência de transmissão e a leitura de identificadores EPC. A análise também considera a seleção de antenas e de diferentes construções de etiquetas, pois o desempenho depende da polarização, da orientação, do material do objeto e da geometria do \textit{inlay}, e não apenas do circuito integrado da tag.

Como aplicação principal, define-se o inventário portátil de bens presentes em uma sala. Pontos fixos de leitura em acessos são discutidos como extensão arquitetural para registrar movimentações, sem confundir a simples detecção de uma tag com a determinação de seu sentido de deslocamento. O protocolo EPC UHF Gen2/ISO/IEC 18000-63 fornece o mecanismo anticolisão da interface aérea; ao software embarcado cabe controlar as rodadas de inventário, receber os quadros, eliminar duplicatas e registrar métricas.

A integração básica e a leitura individual foram obtidas, mas a validação quantitativa de alcance, estabilidade e leitura simultânea ainda requer ensaios controlados. Portanto, o trabalho não afirma a conclusão de um sistema patrimonial completo: estabelece e avalia criticamente o protótipo de leitura que poderá servir de base para essa aplicação.

\textbf{Palavras-chave:} RFID UHF. Inventário patrimonial. YPD-R200. Etiquetas passivas. Anticolisão.

\end{resumo}
